\begin{frame}{XOR의 부등식 모순 — 직접 세워보자}
  \small
  직선분리 가능하다는 말 = 어떤 $(w_1, w_2, b)$가 다음을 \textbf{모두} 만족:
  \[
  w_2 + b > 0\;\; (0,1) \qquad
  w_1 + b > 0\;\; (1,0) \qquad
  b \le 0\;\; (0,0) \qquad
  w_1 + w_2 + b \le 0\;\; (1,1)
  \]
  \begin{enumerate}
    \item 첫 두 식을 더하면 $w_1 + w_2 + 2b > 0$
    \item 셋째 식($b \le 0$)과 결합 $\Rightarrow$ $w_1 + w_2 + b > -b \ge 0$
    \item 넷째 식은 $w_1 + w_2 + b \le 0$이라 주장 — \textbf{모순}
  \end{enumerate}
  \begin{block}{결론}
    부등식 네 개는 어떤 $(w_1, w_2, b)$로도 동시에 성립 불가.
    퍼셉트론에 필요한 것은 \kb{비선형 결정경계}.
  \end{block}
\end{frame}
